% ============================================================================
\clearpage
\section*{v4.10.0-r02 Source-Only: Asymptotics of the Admixture at the Big Bang}
\addcontentsline{toc}{section}{v4.10.0-r02 Source-Only: Asymptotics of the Admixture at the Big Bang}
\label{sec:v4100-r02-admixture-asymptotics}
% ============================================================================

\paragraph{Purpose.}
Under the domain choice G2 the big bang is a transition from the tachyon reading to the matter reading, and a mode of the tachyon
reading picks up an admixture \(C\) of positive frequency after the bang that falls off for short waves (r85.6, \S\ref{sec:v490-r85-masses-cpt}).
Layer v4.9.0-r86 found the fall-off close to \(e^{-\pi k\Delta}\) with an unexplained power prefactor (\S\ref{sec:v490-r86-wolfram-checks}).
This layer records three pre-registered tests on that question (H91--H93).  The last one derives the asymptotics in closed form.  All three have
formal defects in their validity conditions, which are disclosed.  The notes are
\texttt{notes/\allowbreak research/\allowbreak v4\_10\_0\_r02/}.  Earlier layers are not rewritten.

\noindent\textbf{Source-only status.}
This v4.10.0-r02 layer is source-only.  It does not build a PDF, does not assign a new DOI,
and does not authorize a public release.  The v4.10.0 Version DOI
\href{https://doi.org/10.5281/zenodo.23991357}{10.5281/zenodo.23991357} and the Concept DOI
\href{https://doi.org/10.5281/zenodo.19210728}{10.5281/zenodo.19210728} are unchanged by it.

\subsection*{v4.10.0-r02 results}
\label{subsec:v4100-r02-results}

Throughout: \((w\chi')'+k^2w\chi=0\), \(w=e^{Ls}\), \(s=\tfrac12(1+\tanh(t/\Delta))\), \(L=\ln(a/b)\mp i\pi\) for the two senses of rotation, \(a,b\) the
matter weights of r84, \(\Delta=2\); \(C\) is the coefficient of \(e^{-ikt}\) after the bang for \(e^{ikt}\) incoming before it.

\paragraph{Theoremlet v4.10.0-r02.1 (H91: the fall-off up to \(k\Delta=60\)).}
Pre-registered (commit \texttt{a49ed7e}, partly not blind).  Integration along \(t=x+ih\) below the singularity; two contours agree to
\(10^{-10}\) up to \(k\Delta=60\).
\begin{enumerate}[leftmargin=2em]
\item[(a)] Invalid by its own criterion, which was a planning error: agreement with the H87 values to \(10^{-5}\) was required, but H87 had recorded
four digits only.  All deviations, up to \(3.2\times10^{-4}\), lie within that rounding.  By the pre-registration parts (b)--(d) are then invalid.  As
formulated they would also have missed their criteria.
\item[(b)--(d)] A fit with \(\sqrt{k\Delta}\) and \(\ln k\Delta\) gives a slope \(3.049\), \(3\%\) below \(\pi\).  The \(\sqrt{k\Delta}\) term improves the fit only by a
factor \(6.5\).  The slope between \(50\) and \(60\) is \(3.021\).
\end{enumerate}
\emph{Noticed afterwards.}  The local slope is wavy and rises slowly (\(3.026\) at \(60\)).  \(|C|\) is the same for both senses of rotation, to
\(10^{-10}\); the asymmetry of r85.6(d) concerns only the other coefficient.

\paragraph{Theoremlet v4.10.0-r02.2 (H92: the prefactor is the Stokes function of the local problem).}
Pre-registered (commit \texttt{c347a92}, partly not blind).  Near \(t_0=i\pi\Delta/2\), \(s\approx\tfrac12+\Delta/(2\tau)\) with \(\tau=t-t_0\), so \(w\) has an
essential singularity.  In \(\sigma=k\tau\) the local equation is \(y''-(\lambda/\sigma^2)y'+y=0\), with the single parameter \(\lambda=Lk\Delta/2\).  Its
reflection \(r(\lambda)\), continued below \(\sigma=0\), gives \(C\approx r(\lambda)e^{2ikt_0}=r(\lambda)e^{-\pi k\Delta}\).
\begin{enumerate}[leftmargin=2em]
\item[(b)] Confirmed: \(\ln|C|+\pi k\Delta-\ln|r|\) is smooth; the fit residual on \([20,60]\) is \(4.1\times10^{-3}\).  The waves sit in \(r\).
\item[(c)] Confirmed: \(|r(\bar\lambda)|=|r(\lambda)|\) to \(8\times10^{-11}\).
\item[(d)] Confirmed: \(r(\lambda)/(2\pi\lambda)=1-1.6\times10^{-8}\) at \(|\lambda|=10^{-3}\), the Born limit.
\item[(a)] Missed: \(\ln|C|+\pi k\Delta-\ln|r|\to\approx-0.62\), not \(0\).
\end{enumerate}
\emph{Disclosure.}  In run 1 the accuracy control (a contour at \(1.5h\)) lost digits (\(3.7\)--\(5.5\times10^{-6}\), above \(10^{-6}\)).  In run 2 the control
used \(0.7h\) (\(\le6\times10^{-10}\)); the values of \(r\) are identical in both runs.

\paragraph{Theoremlet v4.10.0-r02.3 (H93: the asymptotics derived).}
Pre-registered with the derivation (commit \texttt{a4b1d17}, partly not blind).
\begin{enumerate}[leftmargin=2em]
\item[(i)] With \(\chi=w^{-1/2}u\), \(u''+Qu=0\) and \(w\) going from \(1\) to \(e^L\): \(C=e^{-L/2}C_u\).  This is the factor missing in r02.2(a).
\item[(ii)] For large \(\lambda\) the local \(Q=1-\lambda^2/(4\sigma^4)-\lambda/\sigma^3\) has four turning points.  The over-barrier formula
\(r\approx-i\,e^{2iF}\) at the turning point \(\eta=-i\) (in \(\sigma=\sqrt{\lambda/2}\,\eta\)) gives \(\ln|r|=c\sqrt{k\Delta}\) with
\[
c=\frac{2\sqrt\pi\,\Gamma(3/4)}{\Gamma(1/4)}\,\operatorname{Re}\sqrt L=1.809205,
\]
\[
1+\int_1^\infty\bigl(1-\sqrt{1-\eta^{-4}}\bigr)d\eta=\frac{2\sqrt\pi\,\Gamma(3/4)}{\Gamma(1/4)}=1.19814.
\]
The term \(-\lambda/\sigma^3\) contributes only the phase \(-\pi/4\), so no additive constant appears.
\item[(iii)] Together:
\[
\ln|C|=-\pi k\Delta+c\sqrt{k\Delta}-\tfrac12\ln(a/b)+O\bigl((k\Delta)^{-1/2}\bigr).
\]
\end{enumerate}
Numbers (run 2): slope of \(\ln|r|\) between \(k\Delta=300\) and \(600\) is \(1.808965\) (\(-0.013\%\), criterion \(0.3\%\)); \(\ln|r|-c\sqrt{k\Delta}=0.0057,\ 0.0040\)
(criterion \(\le0.05\), falling); the full problem at \(k\Delta=80,100,120\) deviates from (iii) by \(-0.0028,\ -0.0085,\ -0.0045\) (criterion \(0.02\)).  Both
senses of rotation give the same numbers.
\emph{Disclosure.}  The validity condition (accuracy controls \(\le10^{-6}\)) is not met in one place.  For the sense \(+\) at \(k\Delta=600\) the control is
\(1.4\times10^{-5}\) in run 2 (contour \(1.3h\) instead of \(0.7h\), changed after run 1, where the \(0.7h\) contour lost digits up to \(1.0\)).  Formally the
test is therefore invalid.  For the sense \(-\) it is valid in full and confirmed, and every content criterion is met with a margin of a factor
\(2\) to \(20\).

\paragraph{Proof and measurement sources.}
Scripts \texttt{k206\_\allowbreak abfall\_\allowbreak wurzelterm.py}, \texttt{k207\_\allowbreak lokales\_\allowbreak problem.py},
\texttt{k208\_\allowbreak wkb\_\allowbreak pruefung.py} with their outputs, and the notes in \texttt{notes/\allowbreak research/\allowbreak v4\_10\_0\_r02/};
pre-registrations and results H91--H93 in \texttt{notes/\allowbreak research/\allowbreak predictions/}.  \(\square\)

\paragraph{Consequence.}
The admixture at the big bang now has a closed asymptotic form.  The exponent \(\pi\) comes from the position of the singularity of \(\tanh\), the
\(\sqrt{k\Delta}\) term from the essential singularity of the weight \(e^{Ls}\) there, and the constant from the adiabatic amplitude between the
weights before and after the bang.  This explains the slope below \(\pi\) of r86.1(b), the waves (from the subdominant turning point, not checked
separately) and the missing factor of r02.2(a).  The cut-off scale stays about \(1/(\pi\Delta)\), shifted only by a factor growing like
\(e^{1.8\sqrt{k\Delta}}\); \(\Delta\) remains free.  OP5 is not closed.

\subsection*{v4.10.0-r02 status in the three frames}
\label{subsec:v4100-r02-frames}

\begin{itemize}[leftmargin=2em]
\item The transition joins frame T (before the bang) and frame M (after it) under G2, with the weights \(a\), \(b\) of frame M.  All results concern
this M--T transition.  A frame-free statement is not made.
\item Photon frame P: with photon weights \(a=b\) one has \(L=\mp i\pi\), the constant \(-\tfrac12\ln(a/b)\) vanishes and \(c\) would be
\(1.19814\sqrt{\pi/2}=1.502\).  This is a consequence noticed afterwards and was not computed.  In the sharp limit the photon weights give no
admixture (r85.5), unchanged.
\end{itemize}

\subsection*{v4.10.0-r02 register entries}
\label{subsec:v4100-r02-register}

\begin{center}
\small
\begin{tabular}{@{}>{\raggedright\arraybackslash}p{0.34\textwidth}>{\raggedright\arraybackslash}p{0.58\textwidth}@{}}
\toprule
statement & status after v4.10.0-r02\\
\midrule
fall-off up to \(k\Delta=60\) & numbers secured; H91 invalid (planning error)\\
prefactor as local Stokes function & waves, symmetry and Born limit confirmed; normalization missed (H92)\\
closed asymptotics of \(\ln|C|\) & derived; numbers met; formally invalid for sense \(+\), valid and confirmed for sense \(-\) (H93)\\
\(|C|\) for both senses & equal (afterwards, H91; local form H92c)\\
cut-off scale & \(\approx1/(\pi\Delta)\), \(\Delta\) free\\
\bottomrule
\end{tabular}
\end{center}

\subsection*{v4.10.0-r02 ledger and guard update}
\label{subsec:v4100-r02-ledger-guard}

\[
R_{\rm chk}^{v4100\text{-}r02}=(1,1,0)
\]
for (fall-off computed to \(k\Delta=600\); asymptotics derived; one accuracy control above its bound).

The layer spends no guard:
\[
R_{\rm glob}^{(v4100\text{-}r02,+)}=(0,1,1,1,1),
\qquad
\Gamma_{\rm guard}^{v4100\text{-}r02}=(1,1,1,1,1),
\qquad
\beta_4=1/256\ \text{unspent}.
\]
It does not imply exact finite \(C5\) zero closure, exact all-readout terminality,
\(C6/T2\), a global photon-clock/RH theorem, or global all-readout closure.
